%% BEATRIX Innosuisse Innovation Cheque - Antragstext
%% LaTeX Version für professionelle Formatierung
%% Erstellt: Januar 2026

\documentclass[11pt,a4paper]{article}

% Packages
\usepackage[utf8]{inputenc}
\usepackage[T1]{fontenc}
\usepackage[ngerman]{babel}
\usepackage{geometry}
\usepackage{booktabs}
\usepackage{longtable}
\usepackage{array}
\usepackage{xcolor}
\usepackage{hyperref}
\usepackage{enumitem}
\usepackage{titlesec}
\usepackage{fancyhdr}
\usepackage{parskip}
\usepackage{microtype}

% Page geometry
\geometry{
    left=2.5cm,
    right=2.5cm,
    top=2.5cm,
    bottom=2.5cm
}

% Colors
\definecolor{fehrgray}{RGB}{80,80,80}
\definecolor{fehrblue}{RGB}{0,51,102}
\definecolor{limitation}{RGB}{139,69,19}

% Hyperref setup
\hypersetup{
    colorlinks=true,
    linkcolor=fehrblue,
    citecolor=fehrblue,
    urlcolor=fehrblue
}

% Section formatting
\titleformat{\section}
    {\Large\bfseries\color{fehrblue}}
    {\thesection}{1em}{}
\titleformat{\subsection}
    {\large\bfseries\color{fehrgray}}
    {\thesubsection}{1em}{}

% Header/Footer
\pagestyle{fancy}
\fancyhf{}
\fancyhead[L]{\small\textcolor{fehrgray}{BEATRIX -- Innosuisse Innovation Cheque}}
\fancyhead[R]{\small\textcolor{fehrgray}{FehrAdvice \& Partners AG}}
\fancyfoot[C]{\thepage}
\renewcommand{\headrulewidth}{0.4pt}

% Custom commands
\newcommand{\limitation}[1]{\textcolor{limitation}{\textit{Limitation: #1}}}
\newcommand{\evidence}[1]{\textcolor{fehrgray}{[#1]}}

% Table column types
\newcolumntype{L}[1]{>{\raggedright\arraybackslash}p{#1}}
\newcolumntype{C}[1]{>{\centering\arraybackslash}p{#1}}

\begin{document}

% Title
\begin{center}
    {\LARGE\bfseries\color{fehrblue} BEATRIX}\\[0.5em]
    {\large Behavioral Economics AI for Transformation and\\Research in eXperimental Contexts}\\[1.5em]
    {\Large Innosuisse Innovation Cheque -- Antragstext}\\[1em]
    \rule{0.6\textwidth}{0.5pt}\\[1em]
    {\normalsize
    \textbf{Antragsteller:} FehrAdvice \& Partners AG\\
    \textbf{Forschungspartner:} Universität Zürich\\
    \textbf{Datum:} Januar 2026\\
    \textbf{Fördersumme:} CHF 15'000
    }
\end{center}

\vspace{1em}

\begin{center}
\fbox{\parbox{0.9\textwidth}{
\small\textbf{Kommunikationsstile (Version 2.0):}\\[0.5em]
\textbullet\ \textbf{Ernst Fehr} (Verhaltensökonomie): Evidenzbasiert, Hedging, explizite Limitationen\\
\textbullet\ \textbf{Harald Gall} (Software Engineering): MSR-Terminologie, Architektur-Metriken, Software-Evolution\\
\textbullet\ \textbf{Johannes Luger} (Strategic Management): Attention-Based View, selektive Informationsverarbeitung, Gioia-Methodik
}}
\end{center}

\vspace{2em}

\tableofcontents
\newpage

%% =============================================================================
\section{Innovation Potential}
%% =============================================================================

\subsection{Frage a1: What makes the project idea technologically or business-wise novel?}

\textbf{Wissenschaftliche Grundlage und Forschungslücke:}

Die verhaltensökonomische Forschung der letzten fünf Jahrzehnte hat robust dokumentiert, dass menschliches Entscheidungsverhalten systematisch von den Annahmen des \textit{Homo oeconomicus} abweicht. Experimentelle Evidenz zeigt, dass Fairness-Präferenzen \evidence{Fehr \& Schmidt, 1999; Fehr \& Fischbacher, 2002}, Reziprozität \evidence{Fehr \& Gächter, 2000}, Verlustaversion \evidence{Kahneman \& Tversky, 1979} und zeitliche Inkonsistenz \evidence{Laibson, 1997} robuste Determinanten ökonomischen Verhaltens darstellen.

Trotz dieser wissenschaftlichen Erkenntnisse fehlt bisher ein systematischer Ansatz, diese Befunde für praktische Entscheidungsunterstützung nutzbar zu machen. BEATRIX adressiert diese Lücke.

\textbf{Der methodische Kern:}

BEATRIX operationalisiert zwei theoretische Konstrukte, die in der bisherigen Anwendung unterrepräsentiert sind:

\begin{enumerate}[leftmargin=2em]
    \item \textbf{Kontext ($\Psi$):} Meta-Analysen dokumentieren erhebliche Heterogenität in Interventionseffekten, die primär durch Kontextfaktoren erklärt wird \evidence{DellaVigna \& Linos, 2022}. BEATRIX erfasst diese systematisch in 8 Dimensionen.

    \item \textbf{Komplementarität (C):} Die Evidenz zum Motivations-Crowding \evidence{Gneezy \& Rustichini, 2000} zeigt, dass extrinsische Anreize unter bestimmten Bedingungen intrinsische Motivation verdrängen können. BEATRIX modelliert diese Interaktionseffekte explizit.
\end{enumerate}

\textbf{Technologische Umsetzung:}

\begin{table}[h]
\centering
\begin{tabular}{@{}L{3.5cm}L{5cm}L{5cm}@{}}
\toprule
\textbf{Komponente} & \textbf{Funktion} & \textbf{Wissenschaftliche Basis} \\
\midrule
Semantic Calibration & Extraktion von Kontextdimensionen aus unstrukturierten Daten & Ermöglicht $\Psi$-Quantifizierung ohne aufwändige Primärerhebung \\
LLM Monte Carlo & Systematische Variationsanalysen (n > 10'000) & Liefert Wahrscheinlichkeitsverteilungen mit Konfidenzintervallen \\
FEPSDE-Wohlfahrt & Multidimensionale Nutzenfunktion & Basiert auf erweitertem Utility-Framework \\
\bottomrule
\end{tabular}
\end{table}

\limitation{BEATRIX liefert probabilistische Schätzungen, keine deterministischen Vorhersagen. Die Prognosegenauigkeit hängt von der Kontexterfassung und der Übertragbarkeit der zugrundeliegenden Studien ab.}

%% -----------------------------------------------------------------------------
\subsection{Frage a2: How does it differ from existing solutions in the market?}

\textbf{Systematische Abgrenzung:}

\begin{table}[h]
\centering
\begin{tabular}{@{}L{3.5cm}L{4.5cm}L{5cm}@{}}
\toprule
\textbf{Ansatz} & \textbf{Theoretische Basis} & \textbf{Limitation} \\
\midrule
Traditionelle Beratung & Implizites Erfahrungswissen & Anfällig für kognitive Verzerrungen \evidence{Kahneman, 2011} \\
Generative KI (LLMs) & Statistische Sprachmuster & Keine kausale Verhaltensmodellierung \\
Business Intelligence & Deskriptive Statistik & Korrelation $\neq$ Kausalität \\
\textbf{BEATRIX} & Axiomatische Verhaltensökonomik & Basiert auf experimentell validierter Evidenz \\
\bottomrule
\end{tabular}
\end{table}

\textbf{Der fundamentale Unterschied:}
\begin{itemize}[leftmargin=2em]
    \item \textbf{LLM-Ansatz:} $P(\text{Text} \mid \text{Kontext})$ -- Wahrscheinlichkeit eines Textoutputs
    \item \textbf{BEATRIX-Ansatz:} $P(\text{Verhalten} \mid \text{Kontext}, \text{Anreize}, \text{Präferenzen})$ -- Wahrscheinlichkeit eines Verhaltensoutcomes
\end{itemize}

BEATRIX ist auf 1'547 experimentellen Studien kalibriert (312 Labor-Experimente, 285 Feldexperimente, 438 RCTs, 512 Kontextstudien). Die Parameter werden durch Meta-Analysen geschätzt, nicht durch maschinelles Lernen auf Textkorpora.

\limitation{Die Übertragbarkeit experimenteller Befunde auf neue Kontexte (external validity) ist nicht garantiert.}

%% -----------------------------------------------------------------------------
\subsection{Frage a3: What is the innovative element that research will validate?}

Der Forschungspartner (Universität Zürich) adressiert zwei offene wissenschaftliche Fragen:

\textbf{Forschungsfrage 1: Technische Robustheit (Prof. Harald Gall)}

\textit{Hintergrund:} Verhaltensmodelle können bei Kontextänderungen an Gültigkeit verlieren (``Model Drift''). Die Software-Engineering-Forschung zu Mining Software Repositories (MSR) hat Methoden zur Erkennung solcher Instabilitäten entwickelt \evidence{Gall, Jazayeri \& Krajewski, 2003; Hassan, 2008}. Insbesondere die Analyse von Code-Änderungen auf Ebene abstrakter Syntaxbäume (AST) ermöglicht feingranulare Drift-Diagnostik -- ein Ansatz, der durch ChangeDistiller \evidence{Fluri et al., 2007} für Software-Evolution etabliert wurde.

\textit{Forschungsfrage:} Können Techniken aus dem Mining Software Repositories (MSR) -- insbesondere Change-Impact-Analysen und evolutionäre Kopplung \evidence{Zimmermann et al., 2005} -- adaptiert werden, um Model Drift in verhaltensökonomischen Prognosemodellen zu diagnostizieren?

\textit{Erwarteter Beitrag:}
\begin{itemize}[leftmargin=2em]
    \item Diagnose-Protokoll zur Klassifikation von Prognosefehlern (stochastisches Rauschen vs. gradueller Drift vs. Strukturbruch), basierend auf Software-Evolution-Metriken
    \item Architektur-Blueprint für fehlertolerante, modular entkoppelte Systemkomponenten \evidence{Bass, Clements \& Kazman, 2012}
    \item Dokumentierte Schwellenwerte für automatische Alarmierung, kalibriert an historischen Prognosefehler-Verteilungen
\end{itemize}

\textbf{Forschungsfrage 2: Nutzerakzeptanz (Prof. Johannes Luger)}

\textit{Hintergrund:} Die Literatur zur Algorithm Aversion \evidence{Dietvorst et al., 2015} zeigt, dass Menschen algorithmische Empfehlungen oft ablehnen. Aus strategischer Managementperspektive ist zudem relevant, wie Manager algorithmische Inputs in ihre Aufmerksamkeitsallokation integrieren \evidence{Ocasio, 1997} und wie organisationale Strukturen die selektive Informationsverarbeitung beeinflussen \evidence{Luger, Junge \& Mammen, 2023}.

\textit{Forschungsfrage:} Unter welchen Bedingungen akzeptieren Entscheidungsträger KI-gestützte Verhaltensanalysen? Wie interagieren algorithmische Empfehlungen mit bestehenden Entscheidungsheuristiken und organisationalen Aufmerksamkeitsstrukturen?

\textit{Erwarteter Beitrag:}
\begin{itemize}[leftmargin=2em]
    \item Qualitative Studie (n=10-15) zu Akzeptanzdeterminanten, basierend auf Gioia-Methodologie \evidence{Gioia, Corley \& Hamilton, 2013}
    \item Gestaltungsprinzipien für die Mensch-KI-Schnittstelle unter Berücksichtigung kognitiver Kapazitätsgrenzen
    \item Identifikation von Moderatoren (Erfahrung, Domäne, Entscheidungstyp, hierarchische Position)
\end{itemize}

\limitation{Die Validierung dient der Machbarkeitseinschätzung, nicht dem vollständigen Nachweis.}

\newpage
%% =============================================================================
\section{Market Relevance}
%% =============================================================================

\subsection{Frage b1: What concrete problem does the project solve?}

Die empirische Literatur dokumentiert eine erhebliche Diskrepanz zwischen organisationaler Entscheidungsqualität und verfügbarem Verhaltenswissen:

\begin{itemize}[leftmargin=2em]
    \item \textbf{Change Management:} 60-70\% der Transformationsprojekte erreichen ihre Ziele nicht \evidence{Beer \& Nohria, 2000}. Hauptursache: mangelnde Berücksichtigung von Verhaltensreaktionen.

    \item \textbf{Anreizsysteme:} Monetäre Anreize können kontraproduktiv wirken \evidence{Gneezy \& Rustichini, 2000}, werden aber oft ohne Kontextanalyse implementiert.

    \item \textbf{Policy-Interventionen:} Nudging-Effekte variieren erheblich (Median: 1.4pp); die Heterogenität wird durch Kontextfaktoren erklärt \evidence{DellaVigna \& Linos, 2022}.
\end{itemize}

\textbf{Das Kernproblem:} Organisationen haben ausgereifte Systeme für quantitative Daten, aber keine vergleichbaren Instrumente für systematische Verhaltensanalyse.

Aus strategischer Managementperspektive reflektiert dies ein Problem der Aufmerksamkeitsallokation \evidence{Ocasio, 1997}: Organisationale Entscheidungsträger fokussieren systematisch auf quantifizierbare Metriken, während schwerer messbare Verhaltensfaktoren aus dem Aufmerksamkeitsfokus fallen. Die Forschung zur selektiven Informationsverarbeitung \evidence{Luger, Junge \& Mammen, 2023} zeigt, dass strukturelle Faktoren diese Asymmetrie verstärken können.

\limitation{BEATRIX adressiert dieses Problem, kann aber keine Garantie für bessere Outcomes geben.}

%% -----------------------------------------------------------------------------
\subsection{Frage b2: Who are the target customers?}

Die Literatur legt nahe, dass bestimmte Organisationstypen besonders profitieren könnten:

\begin{enumerate}[leftmargin=2em]
    \item \textbf{Beratungsunternehmen:} Traditionell basierend auf Erfahrungswissen. Probleme: Wissenserosion bei Fluktuation \evidence{Hansen et al., 1999}, schwierige Skalierung.

    \item \textbf{Öffentliche Institutionen:} Die Behavioral Insights-Bewegung \evidence{Thaler \& Sunstein, 2008} hat Potential für evidenzbasierte Policy gezeigt.

    \item \textbf{Unternehmen mit Transformationsbedarf:} Veränderungen scheitern häufig an mangelnder Berücksichtigung von Mitarbeiterreaktionen.
\end{enumerate}

\textbf{Pilotierte Anwendungen (vorläufig):}

\begin{table}[h]
\centering
\begin{tabular}{@{}L{3cm}L{4cm}L{3.5cm}L{3cm}@{}}
\toprule
\textbf{Organisation} & \textbf{Anwendung} & \textbf{Ergebnis} & \textbf{Limitation} \\
\midrule
Verkehrsunternehmen & Kohärenz-Analyse & 3/8 Diskrepanzen & Keine Kontrollgruppe \\
Bundesamt & Subvention vs. Nudging & Nudging 40\% effektiver & Nicht feldvalidiert \\
Finanzinstitut & Kampagnentest & 70\% weniger Varianten & Ext. Validität unklar \\
\bottomrule
\end{tabular}
\end{table}

\limitation{Systematische Wirksamkeitsnachweise erfordern kontrollierte Studien.}

%% -----------------------------------------------------------------------------
\subsection{Frage b3: What is the size of the potential market?}

Eine präzise Marktquantifizierung ist mit Unsicherheit behaftet:

\begin{table}[h]
\centering
\begin{tabular}{@{}L{5cm}L{4cm}C{3cm}@{}}
\toprule
\textbf{Segment} & \textbf{Geschätztes Volumen} & \textbf{Konfidenz} \\
\midrule
Management Consulting (global) & USD 300-350 Mrd. & Mittel \\
Management Consulting (Schweiz) & CHF 3-4 Mrd. & Niedrig \\
Relevanter Teilmarkt (10-20\%) & CHF 300-800 Mio. & Niedrig \\
\bottomrule
\end{tabular}
\end{table}

\textbf{Wettbewerbsposition:} FehrAdvice verfügt über wissenschaftliche Nähe zur verhaltensökonomischen Forschung, eine Datenbank mit $\sim$1'500 codierten Studien und etablierte Kundenbeziehungen im DACH-Raum.

\limitation{Ob diese Assets in nachhaltigen Wettbewerbsvorteil übersetzbar sind, hängt von erfolgreicher Validierung ab.}

%% -----------------------------------------------------------------------------
\subsection{Frage b4: What impact can be expected?}

Die folgenden Schätzungen basieren auf Pilotprojekten und Literatur-Extrapolation. Es sind Hypothesen, keine validierten Kausaleffekte:

\begin{table}[h]
\centering
\begin{tabular}{@{}L{4cm}L{3.5cm}L{4cm}C{2cm}@{}}
\toprule
\textbf{Dimension} & \textbf{Baseline} & \textbf{Mit BEATRIX} & \textbf{Unsicherheit} \\
\midrule
Diagnosezeit & 4-6 Wochen & 1-2 Wochen & $\pm$50\% \\
Prognosegenauigkeit & Nicht erfasst & 65-80\% (Pilotdaten) & $\pm$15pp \\
\bottomrule
\end{tabular}
\end{table}

\textbf{Kritische Einschränkungen:}
\begin{itemize}[leftmargin=2em]
    \item Keine Kontrollgruppen in bisherigen Pilotprojekten
    \item Selbstselektion der Kunden möglich
    \item Hawthorne-Effekt nicht ausschliessbar
\end{itemize}

\newpage
%% =============================================================================
\section{Feasibility \& Risk Reduction}
%% =============================================================================

\subsection{Frage c1: What aspect of feasibility needs to be tested?}

\textbf{1. Technische Machbarkeit (Primärfokus):}

Die zentrale Herausforderung betrifft die Stabilität von Verhaltensmodellen unter dynamischen Bedingungen. Die Software-Engineering-Forschung zu ``concept drift'' und Software-Evolution \evidence{Gama et al., 2014; Lehman, 1980} dokumentiert, dass Modelle an Genauigkeit verlieren können. Diese Erkenntnisse aus dem Mining Software Repositories (MSR) -- einem Forschungsfeld, das Prof. Gall massgeblich mitgeprägt hat \evidence{Gall, Jazayeri \& Krajewski, 2003} -- sind auf Verhaltensmodelle übertragbar.

\begin{table}[h]
\centering
\begin{tabular}{@{}L{3.5cm}L{5cm}L{5cm}@{}}
\toprule
\textbf{Offene Frage} & \textbf{Relevanz} & \textbf{Forschungsansatz} \\
\midrule
Model Drift & Kontextänderungen können Parameter invalidieren & Change-Impact-Analysen \evidence{Zimmermann et al., 2005} \\
Skalierbarkeit & Pilotprojekte nicht automatisch übertragbar & Modulare Architektur \evidence{Bass et al., 2012} \\
Interoperabilität & Integration in IT-Landschaften & API-Design mit SE-Metriken \\
\bottomrule
\end{tabular}
\end{table}

\textbf{2. Wissenschaftliche Machbarkeit:}

Die theoretische Basis ist etabliert \evidence{Kahneman \& Tversky, 1979; Fehr \& Schmidt, 1999}. Die Übertragbarkeit auf neue Kontexte bleibt offen \evidence{Levitt \& List, 2007}.

\textbf{3. Nutzerakzeptanz:}

Die Literatur zu Algorithm Aversion/Appreciation zeigt kontextabhängige Akzeptanz \evidence{Dietvorst et al., 2015; Logg et al., 2019}.

%% -----------------------------------------------------------------------------
\subsection{Frage c2: What proof-of-concept can be achieved with CHF 15'000?}

\textbf{Erreichbar:}

\begin{table}[h]
\centering
\begin{tabular}{@{}L{4cm}L{9cm}@{}}
\toprule
\textbf{Deliverable} & \textbf{Methodischer Ansatz} \\
\midrule
Architektur-Assessment & Review durch Prof. Gall nach Software-Engineering-Prinzipien \\
Drift-Diagnose-Konzept & Literaturbasiert + Adaption für Verhaltensmodelle \\
Akzeptanz-Exploration & Qualitative Studie (n=10-15) nach Gioia et al. (2013) \\
Interface-Mockup & Low-fidelity Prototyp mit Nutzerfeedback \\
\bottomrule
\end{tabular}
\end{table}

\textbf{NICHT erreichbar:}
\begin{itemize}[leftmargin=2em]
    \item Vollständige Produktentwicklung
    \item Randomisierte kontrollierte Studien
    \item Umfassende Feldvalidierung
    \item Markteinführung
\end{itemize}

\textbf{Einordnung:} TRL 4 $\rightarrow$ TRL 5-6 (ambitioniert für das Budget). Der Innovation Cheque ermöglicht eine fundierte Machbarkeitseinschätzung.

%% -----------------------------------------------------------------------------
\subsection{Frage c3: How will the research partner's work reduce risks?}

\textbf{Risiko 1: Model Drift (Prof. Harald Gall)}

\begin{itemize}[leftmargin=2em]
    \item \textit{Problem:} Verhaltensmodelle können bei Kontextänderungen an Gültigkeit verlieren. Die MSR-Forschung zeigt, dass evolutionäre Kopplung und Change-Impact-Analysen frühzeitige Warnsignale liefern können \evidence{Zimmermann et al., 2005}.
    \item \textit{Ansatz:} Prof. Galls Expertise in Software Evolution und MSR \evidence{Gall, Jazayeri \& Krajewski, 2003; Hassan, 2008} ermöglicht die Adaption etablierter Methoden -- insbesondere ChangeDistiller-basierte AST-Differenzierung \evidence{Fluri et al., 2007}.
    \item \textit{Output:} Klassifikationsschema für Prognosefehler (stochastisches Rauschen vs. gradueller Drift vs. Strukturbruch); Schwellenwerte für automatische Alarmierung.
\end{itemize}

\limitation{Konzept wird spezifiziert, nicht vollständig implementiert. Übertragbarkeit auf Verhaltensmodelle ist Forschungsfrage.}

\textbf{Risiko 2: Nutzerakzeptanz (Prof. Johannes Luger)}

\begin{itemize}[leftmargin=2em]
    \item \textit{Problem:} Menschen lehnen algorithmische Empfehlungen oft ab \evidence{Dietvorst et al., 2015}. Aus strategischer Managementperspektive ist dies ein Problem der Aufmerksamkeitsallokation \evidence{Ocasio, 1997} und selektiven Informationsverarbeitung \evidence{Luger, Junge \& Mammen, 2023}.
    \item \textit{Ansatz:} Qualitative Studie basierend auf Gioia-Methodologie \evidence{Gioia, Corley \& Hamilton, 2013}. Fokus auf Sensemaking-Prozesse und Aufmerksamkeitsstrukturen.
    \item \textit{Output:} Akzeptanzdeterminanten; Gestaltungsempfehlungen unter Berücksichtigung kognitiver Kapazitätsgrenzen; testbare Moderationshypothesen.
\end{itemize}

\limitation{Qualitative Exploration erlaubt keine kausalen Schlussfolgerungen; dient der Hypothesengenerierung.}

%% -----------------------------------------------------------------------------
\subsection{Frage c4: What would be the next steps if feasibility is proven?}

\textbf{Drei Szenarien:}

\begin{table}[h]
\centering
\begin{tabular}{@{}L{2.5cm}L{5cm}L{3.5cm}C{2.5cm}@{}}
\toprule
\textbf{Szenario} & \textbf{Beschreibung} & \textbf{Folgeprojekt} & \textbf{Unsicherheit} \\
\midrule
A: Positiv & Architektur skalierbar, Akzeptanz vorhanden & CHF 500k-2M Hauptprojekt & Hoch \\
B: Gemischt & Einzelne Aspekte positiv & Fokussierte Fortsetzung & Mittel \\
C: Negativ & Fundamentale Probleme & Neuausrichtung/Einstellung & -- \\
\bottomrule
\end{tabular}
\end{table}

\limitation{Diese Roadmap ist hypothetisch. Zeitrahmen und Budgets sind Schätzungen mit hoher Unsicherheit.}

\newpage
%% =============================================================================
\section{Collaboration with a Research Partner}
%% =============================================================================

\subsection{Frage d1: What scientific expertise is required that we don't have in-house?}

\textbf{Interne Kompetenzen (FehrAdvice):}

\begin{table}[h]
\centering
\begin{tabular}{@{}L{4cm}L{5cm}L{4.5cm}@{}}
\toprule
\textbf{Bereich} & \textbf{Beschreibung} & \textbf{Einschränkung} \\
\midrule
Verhaltensökonomie & Anwendungserfahrung & Keine Grundlagenforschung \\
Literaturdatenbank & $\sim$1'500 codierte Studien & Sekundäranalyse \\
Pilotprojekte & $\sim$20 Anwendungen & Ohne Kontrollgruppen \\
\bottomrule
\end{tabular}
\end{table}

\textbf{Identifizierte Kompetenzlücken:}

\begin{table}[h]
\centering
\begin{tabular}{@{}L{3cm}L{5cm}L{3.5cm}L{2cm}@{}}
\toprule
\textbf{Kompetenz} & \textbf{Begründung} & \textbf{Literatur} & \textbf{Durch} \\
\midrule
Software Evolution \& MSR & Verhaltensmodelle erfordern kontinuierliche Anpassung; MSR bietet Drift-Erkennung & Gall et al. (2003); Lehman (1980) & Prof. Gall \\
Architektur-Design & Skalierbare Systeme mit Qualitätsmetriken & Bass et al. (2012); Chidamber \& Kemerer (1994) & Prof. Gall \\
Strategic Decision Making & Integration algorithmischer Outputs in Entscheidungen & Ocasio (1997); Luger et al. (2023) & Prof. Luger \\
Human-AI Interaction & Aufmerksamkeitsallokation bei algorithmischen Empfehlungen & Dietvorst et al. (2015) & Prof. Luger \\
Qualitative Methoden & Rigorose Exploration nach Gioia & Gioia et al. (2013) & Prof. Luger \\
\bottomrule
\end{tabular}
\end{table}

%% -----------------------------------------------------------------------------
\subsection{Frage d2: Why is the chosen research partner most suitable?}

\textbf{Prof. Harald Gall (Institut für Informatik, Dekan der Wirtschaftswissenschaftlichen Fakultät):}

Prof. Gall ist international anerkannter Experte für Software Evolution und Mining Software Repositories (MSR) mit $>$20'000 Zitationen. Seine Forschung hat das Feld der automatisierten Software-Analyse massgeblich geprägt:

\begin{itemize}[leftmargin=2em]
    \item \textbf{Mining Software Repositories:} Adaption von MSR-Techniken für Drift-Erkennung \evidence{Gall, Jazayeri \& Krajewski, 2003; Hassan, 2008}
    \item \textbf{Change-Impact-Analyse:} ChangeDistiller-Methodik zur feingranularen Änderungserkennung auf AST-Ebene \evidence{Fluri et al., 2007; Zimmermann et al., 2005}
    \item \textbf{Software-Qualitätsmetriken:} Systematische Architektur-Bewertung \evidence{Chidamber \& Kemerer, 1994}
\end{itemize}

\textbf{Prof. Johannes Luger (Institut für Betriebswirtschaft, Evidence-Based Strategic Management):}

Prof. Luger forscht an der Schnittstelle von strategischem Management, Entscheidungsfindung und Organisationstheorie. Seine Arbeiten zur selektiven Informationsverarbeitung sind direkt relevant:

\begin{itemize}[leftmargin=2em]
    \item \textbf{Selective Information Processing:} Wie Manager algorithmische Outputs integrieren \evidence{Luger, Junge \& Mammen, 2023}
    \item \textbf{Attention-Based View:} Aufmerksamkeitsallokation bei konkurrierenden Informationsquellen \evidence{Ocasio, 1997; König et al., 2018}
    \item \textbf{Qualitative Organisationsforschung:} Rigorose Exploration nach Gioia-Methodologie \evidence{Gioia, Corley \& Hamilton, 2013}
\end{itemize}

\textbf{Begründung UZH vs. ETH/EPFL:}

\begin{table}[h]
\centering
\begin{tabular}{@{}L{5cm}C{3cm}C{3cm}@{}}
\toprule
\textbf{Kriterium} & \textbf{UZH} & \textbf{ETH/EPFL} \\
\midrule
Kombination Informatik + BWL & \checkmark & Primär technisch \\
Verhaltensökonomie-Tradition & \checkmark (Fehr-Schule) & Weniger ausgeprägt \\
\bottomrule
\end{tabular}
\end{table}

%% -----------------------------------------------------------------------------
\subsection{Frage d3: What is the specific contribution of the research partner?}

\textbf{Arbeitspaket 1: Architektur-Assessment (Prof. Gall, $\sim$CHF 8'000)}

\begin{itemize}[leftmargin=2em]
    \item Systematische Architektur-Bewertung $\rightarrow$ Stärken-Schwächen-Analyse
    \item Drift-Erkennung Konzept $\rightarrow$ Adaptiert für Verhaltensmodelle
    \item Review-Kompetenz $\rightarrow$ Priorisierte Empfehlungen
\end{itemize}

\textbf{Arbeitspaket 2: Akzeptanz-Exploration (Prof. Luger, $\sim$CHF 7'000)}

\begin{itemize}[leftmargin=2em]
    \item Qualitative Forschungsmethoden $\rightarrow$ Interviewleitfaden
    \item Netzwerk-Zugang $\rightarrow$ 10-15 Interviews
    \item Thematische Analyse $\rightarrow$ Vorläufiges Akzeptanzmodell
\end{itemize}

\limitation{CHF 15'000 ermöglicht konzeptuelle Grundlagen und erste explorative Daten -- keine fertigen Lösungen.}

%% -----------------------------------------------------------------------------
\subsection{Frage d4: How will this collaboration strengthen innovation capacity?}

Die Literatur zu Absorptive Capacity \evidence{Cohen \& Levinthal, 1990} legt nahe, dass Unternehmen von Forschungskooperationen profitieren können:

\begin{table}[h]
\centering
\begin{tabular}{@{}L{4.5cm}L{5.5cm}C{3cm}@{}}
\toprule
\textbf{Dimension} & \textbf{Erwarteter Effekt} & \textbf{Unsicherheit} \\
\midrule
Technisches Know-how & Transfer von Software-Engineering-Methoden & Mittel \\
Methodisches Repertoire & Erweiterung um qualitative Forschungsmethoden & Mittel \\
Wissenschaftliche Legitimität & Potenzielle Co-Publikationen & Hoch \\
Netzwerkeffekte & Zugang zu akademischer Community & Mittel \\
\bottomrule
\end{tabular}
\end{table}

\limitation{Diese Effekte sind nicht garantiert. University-Industry Collaborations scheitern häufig an Transferbarrieren \evidence{Bruneel et al., 2010}.}

%% -----------------------------------------------------------------------------
\subsection{Frage d5: Could this lead to a longer-term partnership?}

\textbf{Szenarien:}

\begin{table}[h]
\centering
\begin{tabular}{@{}L{2.5cm}L{5cm}L{4cm}C{2cm}@{}}
\toprule
\textbf{Szenario} & \textbf{Folgeprojekt} & \textbf{Voraussetzung} & \textbf{Unsicherheit} \\
\midrule
A: Positiv & Hauptprojekt CHF 500k-2M & Architektur skalierbar & Hoch \\
B: Gemischt & Fokussierte Fortsetzung & Schwachstellen adressierbar & Mittel \\
C: Negativ & Keine Fortsetzung & -- & -- \\
\bottomrule
\end{tabular}
\end{table}

\textbf{Erfolgsfaktoren für langfristige Partnerschaft} \evidence{Perkmann et al., 2013}:

\begin{itemize}[leftmargin=2em]
    \item Komplementäre Kompetenzen: \textit{Vorhanden} (Positiv)
    \item Klare Zieldefinition: \textit{Definiert} (Positiv)
    \item Persönliche Beziehungen: \textit{Zu entwickeln} (Offen)
\end{itemize}

\limitation{Eine langfristige Partnerschaft ist möglich, aber nicht garantiert. Der Innovation Cheque dient auch der Exploration, ob eine solche Partnerschaft sinnvoll ist.}

\newpage
%% =============================================================================
\section{Impact for Switzerland}
%% =============================================================================

\subsection{Frage e1: How does the project strengthen Swiss industry competitiveness?}

Die Frage nach nationalem Wettbewerbsvorteil erfordert kritische Einordnung: Nationale Kompetenzcluster sind historisch gewachsen und nicht deterministisch planbar \evidence{Porter, 1990; Feldman \& Kogler, 2010}.

\textbf{Know-how-Aufbau:}

\begin{table}[h]
\centering
\begin{tabular}{@{}L{3cm}L{4cm}L{3cm}L{3.5cm}@{}}
\toprule
\textbf{Dimension} & \textbf{Hypothese} & \textbf{Evidenzbasis} & \textbf{Limitation} \\
\midrule
Methodische Kompetenz & Nische Verhaltensökonomie + AI & Fehr \& Schmidt, 1999 & Parallele Entwicklungen \\
Humankapital & University-Industry-Spillovers & Perkmann et al., 2013 & Effektgrössen variabel \\
Absorptive Capacity & Externes Wissen nutzen & Cohen \& Levinthal, 1990 & Attribution schwierig \\
\bottomrule
\end{tabular}
\end{table}

\textbf{Beschäftigung (spekulative Schätzung):}

\begin{table}[h]
\centering
\begin{tabular}{@{}L{3.5cm}L{5cm}C{3cm}C{2cm}@{}}
\toprule
\textbf{Szenario} & \textbf{Annahme} & \textbf{Stellen (5 J.)} & \textbf{Konfidenz} \\
\midrule
Optimistisch & Hohe Marktakzeptanz & 30-50 & Niedrig \\
Realistisch & Moderate Adoption & 10-20 & Mittel \\
Konservativ & Langsame Diffusion & 3-8 & Mittel \\
\bottomrule
\end{tabular}
\end{table}

\limitation{Diese Schätzungen basieren auf Analogien und sind spekulativ.}

%% -----------------------------------------------------------------------------
\subsection{Frage e2: How does it support SME-university collaboration?}

Das Projekt entspricht strukturell einem typischen University-Industry-Linkage-Format \evidence{Perkmann \& Walsh, 2007; D'Este \& Patel, 2007}:

\begin{table}[h]
\centering
\begin{tabular}{@{}L{3cm}L{5cm}L{5.5cm}@{}}
\toprule
\textbf{Richtung} & \textbf{Mechanismus} & \textbf{Limitation} \\
\midrule
UZH $\rightarrow$ FehrAdvice & Methodisches Know-how & Transfer taciten Wissens schwierig \\
FehrAdvice $\rightarrow$ UZH & Praxisdaten für Forschung & Datenschutz-Restriktionen \\
\bottomrule
\end{tabular}
\end{table}

\limitation{Die Literatur warnt vor überzogenen Erwartungen an UIL \evidence{Bozeman et al., 2013}. Ob effektiver Wissenstransfer stattfindet, kann erst ex post beurteilt werden.}

%% -----------------------------------------------------------------------------
\subsection{Frage e3: What are the potential societal or economic benefits?}

\textbf{Methodische Vorbemerkung:} Die Quantifizierung gesellschaftlicher Effekte ist methodisch problematisch \evidence{Salter \& Martin, 2001}.

\begin{table}[h]
\centering
\begin{tabular}{@{}L{3.5cm}L{3.5cm}L{3.5cm}C{2.5cm}@{}}
\toprule
\textbf{Bereich} & \textbf{Evidenzbasis} & \textbf{Effektgrösse} & \textbf{Konfidenz} \\
\midrule
Organisationale Entscheidungen & Kahneman et al. (2011) & 10-30\% Varianzreduktion & Mittel \\
Public Policy & DellaVigna \& Linos (2022) & d = 0.05-0.15 (Feld) & Mittel \\
\bottomrule
\end{tabular}
\end{table}

\textbf{Wichtige Limitationen:}
\begin{itemize}[leftmargin=2em]
    \item Attribution des kausalen Beitrags eines Projekts schwierig
    \item Laboreffekte übertragen sich oft nicht 1:1 (Feldeffekte 60-90\% kleiner)
    \item Verhaltensinterventionen wirken heterogen \evidence{Allcott, 2015}
\end{itemize}

\limitation{Gesellschaftlicher Nutzen ist plausibel, aber nicht quantifizierbar. Überzogene Behauptungen würden wissenschaftlicher Redlichkeit widersprechen.}

%% -----------------------------------------------------------------------------
\subsection{Frage e4: Could the project help position Switzerland as a leader?}

\textbf{Kritische Analyse:} Führerschaft in Innovationsfeldern ist historisch gewachsen, von kritischer Masse abhängig und pfadabhängig \evidence{Porter, 1990; Feldman, 2000}.

\textbf{Bestandsaufnahme:}

\begin{table}[h]
\centering
\begin{tabular}{@{}L{4.5cm}L{4cm}L{4cm}@{}}
\toprule
\textbf{Dimension} & \textbf{Status} & \textbf{International} \\
\midrule
Verhaltensökonomische Forschung & Stark & Top 5 weltweit \\
AI/ML Forschung & Stark & Top 10 weltweit \\
Kombination Behavioral + AI & Emergent & Unklar \\
Industrielle Anwendung & Begrenzt & Keine kritische Masse \\
\bottomrule
\end{tabular}
\end{table}

\textbf{Realistische Einschätzung:}

\begin{table}[h]
\centering
\begin{tabular}{@{}L{6cm}C{4cm}L{3.5cm}@{}}
\toprule
\textbf{Szenario} & \textbf{Wahrscheinlichkeit} & \textbf{Konditionen} \\
\midrule
Nischenführer DACH & Mittel-Hoch & Erfolgreiche Skalierung \\
Beitrag zu internationalem Netzwerk & Hoch & Kooperationen \\
Globaler Marktführer & Niedrig & Massive Investitionen \\
\bottomrule
\end{tabular}
\end{table}

\limitation{Das Projekt kann einen bescheidenen Beitrag zur Schweizer Positionierung leisten. Die Behauptung ``globaler Führerschaft'' wäre unbegründet und würde den Standards wissenschaftlicher Bescheidenheit widersprechen.}

\newpage
%% =============================================================================
\section{Zusammenfassung}
%% =============================================================================

\begin{center}
\fbox{\parbox{0.95\textwidth}{
\textbf{Abstract für Innolink-Portal:}

\vspace{0.5em}

BEATRIX (Behavioral Economics AI for Transformation and Research in eXperimental contexts) adressiert eine dokumentierte Forschungslücke: Die verhaltensökonomische Literatur zeigt, dass Kontext ($\Psi$) und Komplementarität (C) entscheidend für Interventionserfolge sind \evidence{Fehr \& Schmidt, 1999; DellaVigna \& Linos, 2022}, aber systematische Anwendungstools fehlen.

\vspace{0.5em}

Der Innovation Cheque finanziert eine Vorstudie mit der Universität Zürich zu: (1) Model-Drift-Erkennung in Verhaltensmodellen (Prof. Gall) und (2) Nutzerakzeptanz von KI-gestützter Verhaltensanalyse (Prof. Luger).

\vspace{0.5em}

\textbf{Erwartete Outputs:} Architektur-Assessment, Drift-Diagnose-Konzept, qualitative Akzeptanzstudie (n=10-15), Interface-Mockup.

\vspace{0.5em}

Die Ergebnisse dienen der Go/No-Go-Entscheidung für ein grösseres Innosuisse-Innovationsprojekt. Positive wie negative Ergebnisse haben wissenschaftlichen Wert.
}}
\end{center}

\vspace{1em}

\textbf{Übersicht aller 20 Fragen:}

\begin{table}[h]
\centering
\begin{tabular}{@{}L{1cm}L{6cm}C{2cm}C{4cm}@{}}
\toprule
\textbf{Nr.} & \textbf{Rubrik} & \textbf{Fragen} & \textbf{Stil} \\
\midrule
a) & Innovation Potential & 3 & Evidenzbasiert mit Limitationen \\
b) & Market Relevance & 4 & Evidenzbasiert mit Limitationen \\
c) & Feasibility \& Risk Reduction & 4 & Evidenzbasiert mit Limitationen \\
d) & Collaboration with Research Partner & 5 & Evidenzbasiert mit Limitationen \\
e) & Impact for Switzerland & 4 & Evidenzbasiert mit Limitationen \\
\midrule
& \textbf{Total} & \textbf{20} & \\
\bottomrule
\end{tabular}
\end{table}

\vspace{2em}

\begin{center}
\rule{0.4\textwidth}{0.5pt}\\[0.5em]
{\small\textit{Dokument erstellt: Januar 2026 | Version 1.0 | Für Innosuisse Innovation Cheque Antrag BEATRIX}}\\[0.5em]
{\small FehrAdvice \& Partners AG | Universität Zürich}
\end{center}

\end{document}
